\centerline{\bf \Huge Логика всего и вся II}

\begin{flushright}
        \scriptsize{Единственный счастливый предмет, который вам понадобится, уже находится внутри — УВЕРЕННОСТЬ\\ Мидорима Синтаро}
\end{flushright}



\problem{\textbf{1}} 
Есть три комнаты, на двери каждой из них — табличка. А написано на табличках вот что:

На первой: «В этой комнате сидит дракон».

На второй: «В этой комнате — принцесса».

На третьей: «Дракон сидит во второй комнате».

Написанное на этих табличках может оказаться правдой, а может и нет; известно, однако, что только на одной из них — правда. А еще мы знаем, что принцесса — лишь в одной из комнат, а в двух других — драконы. Так где же сидит принцесса?\\


\problem{\textbf{2}} На острове каждый житель либо рыцарь (всегда говорит правду), либо
лжец (всегда лжёт), либо обычный человек (может и говорить правду, и лгать). Жители этого
острова А и В сказали следующее. А: «B — рыцарь». В: «А — не рыцарь». Докажите, что по
крайней мере один из них говорит правду, но это не рыцарь.\\



\problem{\textbf{3}} На острове Контрастов живут и рыцари, и лжецы.
Рыцари всегда говорят правду, лжецы всегда лгут. Некоторые жители заявили, что на острове
чётное число рыцарей, а остальные заявили, что на острове нечётное число лжецов. Может ли
число жителей острова быть нечётным?\\


\problem{\textbf{4}}  В самолёте летят жители города лжецов и жители города рыцарей. Рыцари всегда говорят правду, а лжецы всегда обманывают. Все пассажиры сели
в ряды по 4 человека, и бортпроводник задал каждому пассажиру один и тот же вопрос. «Верно
ли, что в вашем ряду столько же Ваших земляков, сколько жителей другого города?» Прозвучало ровно 70 утвердительных ответов. Сколько лжецов летит в самолёте? Человек считается
своим собственным земляком.\\

\newpage


\problem{\textbf{5}} За круглым столом сидело 13 гостей. Среди них есть рыцари
(всегда говорят правду), лжецы (всегда лгут) и марсиане. Про марсиан известно, что правду они
говорят только марсианам, а всем остальным лгут. Известно, что каждый сидящий за столом
сказал своему соседу справа: «Ты — лжец». Сколько рыцарей сидело за столом? Найдите все
возможные варианты ответа и докажите, что других нет.\\




\problem{\textbf{6}} Каждый из 10 гномов либо всегда говорит правду,
либо всегда лжёт. Известно, что каждый из них любит ровно один сорт мороженого: сливочное, шоколадное или фруктовое. Сначала Белоснежка попросила поднять руки тех, кто любит
сливочное мороженое, и все подняли руки, потом тех, кто любит шоколадное мороженое — и
половина гномов подняли руки, потом тех, кто любит фруктовое мороженое — и руку поднял
только один гном. Сколько среди гномов правдивых?\\



\problem{\textbf{7}} За круглым столом сидят 38 попугаев и Мартышка. Известно,
что каждый из них либо всегда лжёт (таких будем называть «лжецами»), либо всегда говорит правду (таких будем называть «правдивыми»). Мартышка задала каждому попугаю один
и тот же вопрос: «Кем является Ваш сосед справа — правдивым или лжецом?» (опрос шёл
последовательно по кругу). Первые два попугая (справа от Мартышки) ответили: «мой сосед справа — лжец». Следующие два: «мой сосед справа — правдивый», следующие два: «мой
сосед справа — лжец», и так далее. По окончании опроса Мартышка сказала: «Среди нас не менее 9 правдивых». Сколько правдивых было на самом деле?

